\MWChapter{Governance, review, and theory risk}{治理、评审与理论风险}
\label{ch:safety}

\section{Governance boundary / 治理边界}

The current formal evidence is deliberately governed as high-impact theory work. The specification,
RFC, proof manifest, and delivery record distinguish automated proof checking from independent
review and from a human FCP or ADR decision. A DRI can prepare evidence; that fact does not make
the DRI an independent reviewer \citep{cantiluneSpec,cantiluneRFCTwo,cantiluneClosure}.

\begin{figure}[H]
\centering
\begin{tikzpicture}[x=1.18cm,y=0.82cm,font=\scriptsize\sffamily]
  \foreach \x in {1,...,5}{
    \foreach \y in {1,...,5}{
      \pgfmathtruncatemacro{\r}{\x*\y}
      \ifnum\r<6
        \def\cellcolor{MWMint!18}
      \else\ifnum\r<13
        \def\cellcolor{MWAmber!20}
      \else
        \def\cellcolor{MWBlush!22}
      \fi\fi
      \fill[\cellcolor] (\x-0.5,\y-0.5) rectangle (\x+0.5,\y+0.5);
      \draw[white,line width=1pt] (\x-0.5,\y-0.5) rectangle (\x+0.5,\y+0.5);
    }
  }
  \draw[MWLine] (0.5,0.5) rectangle (5.5,5.5);
  \foreach \x in {1,...,5}\node[below=3pt,text=MWMuted] at (\x,0.5) {\x};
  \foreach \y in {1,...,5}\node[left=3pt,text=MWMuted] at (0.5,\y) {\y};
  \node[below=15pt,font=\small\sffamily,text=MWMuted] at (3,0.5) {Likelihood / 可能性};
  \node[rotate=90,above=16pt,font=\small\sffamily,text=MWMuted] at (0.5,3) {Impact / 影响};

  \node[circle,fill=MWBlush,text=white,inner sep=2.4pt] at (4,5) {R1};
  \node[diamond,fill=MWLilac,text=white,inner sep=2.7pt] at (4,4) {R2};
  \node[rectangle,rounded corners=1pt,fill=MWAmber,text=white,inner sep=2.5pt] at (3,4) {R3};
  \node[circle,draw=MWCyan,line width=1pt,fill=white,text=MWInk,inner sep=2.1pt] at (2,3) {R4};
\end{tikzpicture}
\caption{Current pre-FCP theory risk register. R1: claim exceeds proved scope; R2: no independent
review or FCP disposition; R3: product conformance remains uninstantiated; R4: source/evidence
binding or toolchain verification becomes stale. Marker shape supplements color.}
\label{fig:risk-matrix}
\MWSourceNote{Assessment derived from the current formal specification, proof manifest, and
theory-closure delivery record \citep{cantiluneSpec,cantiluneManifest,cantiluneClosure}.}
\end{figure}


\section{Current risk register / 当前风险登记}

\begin{table}[H]
\centering
\caption{Risks that remain material even when the generic proof manifest reports \emph{proved}.}
\label{tab:controls}
\begin{tabularx}{\textwidth}{@{}L{0.23\textwidth}Y L{0.22\textwidth}@{}}
\toprule
\rowcolor{MWSoft!70}
\textbf{Risk} & \textbf{Current control or evidence} & \textbf{Residual condition} \\
\midrule
R1: scope overclaim & This report labels all formal, product, review, and deployment boundaries explicitly & Readers may still infer runtime capability unless every public summary preserves the boundary \\
R2: non-independent review & Manifest and delivery record keep review evidence, reviewer identity, and conflict-of-interest fields explicit & Independent QA-L4 review and FCP disposition are pending \\
R3: product conformance gap & Generic certificate interfaces and reference witnesses are recorded; eight planned packages remain deferred & A concrete package can diverge unless its own rules and certificates are reviewed \\
R4: stale provenance & Verified source, evidence commit, manifest, and integrity gate are recorded & Later source, toolchain, or documentation changes require renewed binding and review \\
\bottomrule
\end{tabularx}
\end{table}

\section{Review decision / 评审决策}

\begin{MWDecision}
\textbf{Research-only status.} The report accepts the generic theory record as \emph{proved /
review-pending}; it does not authorize deployment, production use, or a product safety claim.
Any promotion requires independent review evidence, the applicable FCP and ADR outcome, and
product-specific conformance certificates.
\end{MWDecision}

\section{Responsible communication / 负责任的沟通}

The theory explains how a future orchestration system could make authority, resources, sessions,
and topology changes inspectable. It does not prove that a real deployment will prevent harmful
actions, satisfy a policy regime, protect personal data, or meet a service-level objective.
Those would require implementation, threat modeling, interface tests, operational controls, and
domain-specific human review beyond the scope of this branch.
